%%%%%%%%%%%%
% PACKAGES %
%%%%%%%%%%%%

\usepackage[utf8]{inputenc}
\usepackage[T1]{fontenc}
\usepackage{polynom}
\usepackage{ulem}
\usepackage{textcomp}
\usepackage{media9}
\usepackage{longdivision}
\usepackage{csquotes}
\usepackage{url}
\usepackage{graphicx}
\usepackage{float}
\usepackage{booktabs}
\usepackage{enumitem}
\usepackage{import}
\usepackage{xifthen}
\usepackage{pdfpages}
\usepackage{transparent}
\usepackage{parskip}
\usepackage{emptypage}
\usepackage{subcaption}
\usepackage{multicol}
\usepackage{cmbright}
\usepackage{xcolor}
\usepackage{amsmath}
\usepackage{amsfonts}
\usepackage{mathtools}
\usepackage{amsthm}
\usepackage{amssymb}
\usepackage{mathrsfs}
\usepackage{cancel}
\usepackage{bm}
\usepackage{systeme}
\usepackage{stmaryrd}
\usepackage{siunitx}
\usepackage{mdframed}
\usepackage{etoolbox}
\usepackage{fancyhdr}
\usepackage{todonotes}
\usepackage{tcolorbox}
\usepackage{erewhon, cabin}
\usepackage[english]{babel}
\usepackage{titletoc}
\usepackage{lipsum}
\usepackage[nottoc]{tocbibind}
\usepackage{thmtools}
\usepackage{thm-restate}
\usepackage{multirow}
\usepackage{hyperref}
\usepackage[noabbrev]{cleveref}
\usepackage{pgfplots}
\usepackage{lmodern}
\usepackage{microtype}
\usepackage{upgreek}
\usepackage[misc]{ifsym}
\usepackage{tikz}
\usepackage{tikz-cd}
\usepackage{tkz-euclide}
\usepackage[
    sorting=nyt,
    style=alphabetic
]{biblatex}


%%%%%%%%%%%%%%%%%%%%%%%%%%%%%%%%%%%%
% New commands and Re-New commands %
%%%%%%%%%%%%%%%%%%%%%%%%%%%%%%%%%%%%

\newcommand{\newtodo}[1]{\todo[color=orange]{TODO:                              \hspace{26.5pt}#1.}}
\newcommand{\canwait}[1]{\todo[color=green]{CAN WAIT:                           \hspace{9.5pt}#1.}}
\newcommand{\important}[1]{\todo[color=yellow]{IMPORTANT:                       #1.}}
\newcommand{\urgent}[1]{\todo[color=red]{URGENT:                                \hspace{16.5pt}#1.}}
\newcommand{\done}[1]{\todo[disable]{#1}\addcontentsline{tdo}{todo}{\sout{DONE: \hspace{26.5pt}#1.}}}

\renewcommand{\bf}{\textbf}
\renewcommand{\it}{\textit}
\renewcommand{\rm}{\textrm}
\newcommand{\un}{\underline}

\newcommand\contra{\scalebox{1.5}{$\lightning$}}
\newcommand\correct[2]{\ensuremath{\:}{\color{red}{#1}}\ensuremath{\to }{\color{correct}{#2}}\ensuremath{\:}}
\newcommand\green[1]{{\color{correct}{#1}}}
\newcommand\hr{
    \noindent\rule[0.5ex]{\linewidth}{0.5pt}                        % Horizontal rule
}
\newcommand\hide[1]{}                                               % Hide parts
\setcounter{secnumdepth}{0}

\newcommand\circled[1]{
    \begin{tikzpicture}[baseline=(char.base)]%
        \node[circle,draw,inner sep=1pt] (char) {\textsf{#1}};%
\end{tikzpicture}}
\newcommand\mc[1]{\footnotesize\circled{#1}}
\newcommand{\incfig}[1]{%
    \def\svgwidth{\columnwidth}
    \import{./figures/}{#1.pdf_tex}
}

\let\oldcite\cite
\renewcommand*\cite[1]{\textsubscript{ (\textit{\oldcite{#1}}) }}

\newcommand{\E}{\mathrm{e}}
\newcommand{\C}{\ensuremath{\mathbb{C}}}
\newcommand{\N}{\ensuremath{\mathbb{N}}}
\newcommand{\Q}{\ensuremath{\mathbb{Q}}}
\newcommand{\R}{\ensuremath{\mathbb{R}}}
\newcommand{\Z}{\ensuremath{\mathbb{Z}}}

\renewcommand{\O}{\ensuremath{\mathbb{O}}}

\let\implies\Rightarrow
\let\impliedby\Leftarrow
\let\iff\Leftrightarrow
\let\epsilon\varepsilon

\newcommand{\fullwidthincfig}[2][0.90]{%
    \ifworking{\makebox[0pt][l]{\color{gray}{\scriptsize\textsf{#2}}}}\fi%
    \def\svgwidth{#1\paperwidth}
    \import{./figures/}{#2.pdf_tex}
}
\newcommand{\minifig}[2]{%
    \def\svgwidth{#1}%
    \begingroup%
    \setbox0=\hbox{\import{./figures/}{#2.pdf_tex}}%
    \parbox{\wd0}{\box0}\endgroup%
    \hspace*{0.2cm}
}
\newcommand\charlotte[1]{{\color{blue}{#1}}}

\def\block(#1,#2)#3{\multicolumn{#2}{c}{\multirow{#1}{*}{$ #3 $}}}


%%%%%%%%%%%%%%%%%%%%%%%
% BEGINNING OF MAKEAT %
%%%%%%%%%%%%%%%%%%%%%%%

\makeatletter
\def\thm@space@setup{ 
    \thm@preskip=\parskip \thm@postskip=0pt                         % Fix some spacing
}

\def\testdateparts#1{\dateparts#1\relax}
\def\dateparts#1 #2 #3 #4 #5\relax{
    \marginpar{\small\textsf{\mbox{#1 #2 #3 #5}}}
}

\def\@lesson{}
\def\@lessonabbr{}
\newcommand{\lesson}[4]{
    \ifthenelse{\isempty{#3}}{
        \def\@lesson{Lesson #1}
        \def\@fullLesson{\hfill \small{#2}\hr \@lesson \hfill \small{#4} \\}
        \subsection[\@lesson]{\@fullLesson}
    }{
        \def\@lesson{Lesson #1: #3}
        \def\@fullLesson{\hfill \small{#2}\hr \@lesson \hfill \small{#4} \\}
        \subsection[\@lesson]{\@fullLesson}
    }
    \def\@lessonabbr{Les #1}
}

\def\@lecture{}
\def\@lectureabbr{}
\newcommand{\lecture}[4]{
    \ifthenelse{\isempty{#3}}{
        \def\@lecture{Lecture #1}
        \def\@fullLecture{\hfill \small{#2}\hr \@lecture \hfill \small{#4} \\}
        \subsection[\@lecture]{\@fullLecture}
    }{
        \def\@lecture{Lecture #1: #3}
        \def\@fullLecture{\hfill \small{#2}\hr \@lecture \hfill \small{#4} \\}
        \subsection[\@lecture]{\@fullLecture}
    }
    \def\@lectureabbr{Lec #1}
}

% These are the fancy headers
\mdfsetup{skipabove=1em,skipbelow=0em, innertopmargin=12pt, innerbottommargin=8pt}
\renewcommand{\sectionmark}[1]{\markboth{#1}{}}                     % Set the \leftmark
\fancypagestyle{plain}{                                             % Changes footer and header
    \fancyhf{}

    \fancyhead[R,L]{\hr}
    \fancyhead[C]{\leftmark}                                        % Center
    \fancyhead[RO,LE]{\@lessonabbr}                                 % Right odd,  Left even
    \fancyhead[RE,LO]{\thepage}                                     % Right even, Left odd
    
    \fancyfoot[RE,LO]{\theTitle}
    \fancyfoot[RO,LE]{\theDate}
    \fancyfoot[C]{\theAuthor}
    
    \renewcommand{\headrulewidth}{0.4pt}
    \renewcommand{\footrulewidth}{0.4pt}
}

\def\pld@CF@loop#1+{%
    \ifx\relax#1\else
        \begingroup
          \pld@AccuSetX11%
          \def\pld@frac{{}{}}\let\pld@symbols\@empty\let\pld@vars\@empty
          \pld@false
          #1%
          \let\pld@temp\@empty
          \pld@AccuIfOne{}{\pld@AccuGet\pld@temp
                            \edef\pld@temp{\noexpand\pld@R\pld@temp}}%
           \pld@if \pld@Extend\pld@temp{\expandafter\pld@F\pld@frac}\fi
           \expandafter\pld@CF@loop@\pld@symbols\relax\@empty
           \expandafter\pld@CF@loop@\pld@vars\relax\@empty
           \ifx\@empty\pld@temp
               \def\pld@temp{\pld@R11}%
           \fi
          \global\let\@gtempa\pld@temp
        \endgroup
        \ifx\@empty\@gtempa\else
            \pld@ExtendPoly\pld@tempoly\@gtempa
        \fi
        \expandafter\pld@CF@loop
    \fi}
\def\pld@CMAddToTempoly{%
    \pld@AccuGet\pld@temp\edef\pld@temp{\noexpand\pld@R\pld@temp}%
    \pld@CondenseMonomials\pld@false\pld@symbols
    \ifx\pld@symbols\@empty \else
        \pld@ExtendPoly\pld@temp\pld@symbols
    \fi
    \ifx\pld@temp\@empty \else
        \pld@if
            \expandafter\pld@IfSum\expandafter{\pld@temp}%
                {\expandafter\def\expandafter\pld@temp\expandafter
                    {\expandafter\pld@F\expandafter{\pld@temp}{}}}%
                {}%
        \fi
        \pld@ExtendPoly\pld@tempoly\pld@temp
        \pld@Extend\pld@tempoly{\pld@monom}%
    \fi}
\makeatother


%%%%%%%%%%%%%%%%%
% END OF MAKEAT %
%%%%%%%%%%%%%%%%%


%%%%%%%%%%%%
% Settings %
%%%%%%%%%%%%

\sisetup{locale = FR}                     % Si unitx
\let\svlim\lim\def\lim{\svlim\limits}
\tcbuselibrary{breakable}                 % Make boxes breakable
\pdfsuppresswarningpagegroup=1            % Fix some stuff
\addbibresource{bibliography.bib}
\hypersetup{                              % Make table of contents clickable
    colorlinks,
    linkcolor={black},
    citecolor={black},
    urlcolor={blue!80!black}
}

\let\cleardoublepage\clearpage            % Clear double page
\newlength{\drop}


%%%%%%%%%%%%%%%%%%%%
% New environments %
%%%%%%%%%%%%%%%%%%%%

\makeatother
\mdfsetup{skipabove=1em,skipbelow=0em}

\tcbuselibrary{skins}

% Color definitions

\definecolor{proofcolor}{RGB}{0,0,0}

% Dark orange and Dark Red rgb
\definecolor{theorembordercolor}{RGB}{151, 63, 5}
\definecolor{theorembackgroundcolor}{RGB}{248, 241, 234}

\definecolor{examplebordercolor}{RGB}{0, 110, 184}
\definecolor{examplebackgroundcolor}{RGB}{240, 244, 250}

\definecolor{definitionbordercolor}{RGB}{0, 150, 85}
\definecolor{definitionbackgroundcolor}{RGB}{239, 247, 243}

\definecolor{propertybordercolor}{RGB}{223, 255, 0}
\definecolor{propertybackgroundcolor}{RGB}{246, 250, 220}


\newtheoremstyle{theorem}
    {0pt}{0pt}{\normalfont}{0pt}
    {}{\;}{0.25em}
    {{\sffamily\bfseries\color{theorembordercolor}\thmname{#1}~\thmnumber{\textup{#2}}.}
        \thmnote{\normalfont\color{black}~(#3)}}
\newtheoremstyle{definition}
    {0pt}{0pt}{\normalfont}{0pt}
    {}{\;}{0.25em}
    {{\sffamily\bfseries\color{definitionbordercolor}\thmname{#1}~\thmnumber{\textup{#2}}.}
        \thmnote{\normalfont\color{black}~(#3)}}
\newtheoremstyle{example}
    {0pt}{0pt}{\normalfont}{0pt}
    {}{\;}{0.25em}
    {{\sffamily\bfseries\color{examplebordercolor}\thmname{#1}.}
        \thmnote{\normalfont\color{black}~(#3)}}
\newtheoremstyle{property}
    {0pt}{0pt}{\normalfont}{0pt}
    {}{\;}{0.25em}
    {{\sffamily\bfseries\color{propertybordercolor}\thmname{#1}~\thmnumber{\textup{#2}}.}
        \thmnote{\normalfont\color{black}~(#3)}}

%%%%%%%%%%%%%%%%%%%%%%%%
% Theorem Environments %
%%%%%%%%%%%%%%%%%%%%%%%%

\makeatother
\mdfsetup{skipabove=1em,skipbelow=0em}

\tcbuselibrary{skins}

% Color definitions

\definecolor{proofcolor}{RGB}{0,0,0}

% Dark orange and Dark Red rgb
\definecolor{theorembordercolor}{RGB}{151, 63, 5}
\definecolor{theorembackgroundcolor}{RGB}{248, 241, 234}

\definecolor{examplebordercolor}{RGB}{0, 110, 184}
\definecolor{examplebackgroundcolor}{RGB}{240, 244, 250}

\definecolor{definitionbordercolor}{RGB}{0, 150, 85}
\definecolor{definitionbackgroundcolor}{RGB}{239, 247, 243}

\definecolor{propertybordercolor}{RGB}{223, 255, 0}
\definecolor{propertybackgroundcolor}{RGB}{246, 250, 220}

\definecolor{solutionbordercolor}{RGB}{26, 26, 26}


\newtheoremstyle{theorem}
    {0pt}{0pt}{\normalfont}{0pt}
    {}{\;}{0.25em}
    {{\sffamily\bfseries\color{theorembordercolor}\thmname{#1}~\thmnumber{\textup{#2}}.}
        \thmnote{\normalfont\color{black}~(#3)}}
\newtheoremstyle{definition}
    {0pt}{0pt}{\normalfont}{0pt}
    {}{\;}{0.25em}
    {{\sffamily\bfseries\color{definitionbordercolor}\thmname{#1}~\thmnumber{\textup{#2}}.}
        \thmnote{\normalfont\color{black}~(#3)}}
\newtheoremstyle{example}
    {0pt}{0pt}{\normalfont}{0pt}
    {}{\;}{0.25em}
    {{\sffamily\bfseries\color{examplebordercolor}\thmname{#1}.}
        \thmnote{\normalfont\color{black}~(#3)}}
\newtheoremstyle{property}
    {0pt}{0pt}{\normalfont}{0pt} {}{\;}{0.25em}
    {{\sffamily\bfseries\color{propertybordercolor}\thmname{#1}~\thmnumber{\textup{#2}}.}
        \thmnote{\normalfont\color{black}~(#3)}}
\newtheoremstyle{solution}
    {0pt}{0pt}{\normalfont}{0pt} {}{\;}{0.25em}
    {{\sffamily\bfseries\color{solutionbordercolor}\thmname{#1}.}
        \thmnote{\normalfont\color{black}~(#3)}}

%%%%%%%%%%%%%%%%%%%%%%%%
% Theorem Environments %
%%%%%%%%%%%%%%%%%%%%%%%%

\theoremstyle{theorem}

\newtheorem{theorem}{Theorem}
\newtheorem{postulate}{Postulate}
\newtheorem{conjecture}{Conjecture}
\newtheorem{corollary}{Corollary}
\newtheorem{lemma}{Lemma}
\newtheorem{conclusion}{Conclusion}

\tcolorboxenvironment{theorem}{
    enhanced jigsaw, pad at break*=1mm, breakable,
    left=4mm, right=4mm, top=1mm, bottom=1mm,
    colback=theorembackgroundcolor, boxrule=0pt, frame hidden,
    borderline west={0.5mm}{0mm}{theorembordercolor}, arc=.5mm
}
\tcolorboxenvironment{postulate}{
    enhanced jigsaw, pad at break*=1mm, breakable,
    left=4mm, right=4mm, top=1mm, bottom=1mm,
    colback=theorembackgroundcolor, boxrule=0pt, frame hidden,
    borderline west={0.5mm}{0mm}{theorembordercolor}, arc=.5mm
}
\tcolorboxenvironment{conjecture}{
    enhanced jigsaw, pad at break*=1mm, breakable,
    left=4mm, right=4mm, top=1mm, bottom=1mm,
    colback=theorembackgroundcolor, boxrule=0pt, frame hidden,
    borderline west={0.5mm}{0mm}{theorembordercolor}, arc=.5mm
}
\tcolorboxenvironment{corollary}{
    enhanced jigsaw, pad at break*=1mm, breakable,
    left=4mm, right=4mm, top=1mm, bottom=1mm,
    colback=theorembackgroundcolor, boxrule=0pt, frame hidden,
    borderline west={0.5mm}{0mm}{theorembordercolor}, arc=.5mm
}
\tcolorboxenvironment{lemma}{
    enhanced jigsaw, pad at break*=1mm, breakable,
    left=4mm, right=4mm, top=1mm, bottom=1mm,
    colback=theorembackgroundcolor, boxrule=0pt, frame hidden,
    borderline west={0.5mm}{0mm}{theorembordercolor}, arc=.5mm
}
\tcolorboxenvironment{conclusion}{
    enhanced jigsaw, pad at break*=1mm, breakable,
    left=4mm, right=4mm, top=1mm, bottom=1mm,
    colback=theorembackgroundcolor, boxrule=0pt, frame hidden,
    borderline west={0.5mm}{0mm}{theorembordercolor}, arc=.5mm
}

%%%%%%%%%%%%%%%%%%%%%%%%%%%
% Definition Environments %
%%%%%%%%%%%%%%%%%%%%%%%%%%%

\theoremstyle{definition}
\newtheorem{definition}{Definition}

\tcolorboxenvironment{definition}{
    enhanced jigsaw, pad at break*=1mm, breakable,
    left=4mm, right=4mm, top=1mm, bottom=1mm,
    colback=definitionbackgroundcolor, boxrule=0pt, frame hidden,
    borderline west={0.5mm}{0mm}{definitionbordercolor}, arc=.5mm
}


%%%%%%%%%%%%%%%%%%%%%%%%
% Example Environments %
%%%%%%%%%%%%%%%%%%%%%%%%

\theoremstyle{example}
\newtheorem*{example}{Example}
\newtheorem*{remark}{Remark}

\tcolorboxenvironment{example}{
    enhanced jigsaw, pad at break*=1mm, breakable,
    left=4mm, right=4mm, top=1mm, bottom=1mm,
    colback=examplebackgroundcolor, boxrule=0pt, frame hidden,
    borderline west={0.5mm}{0mm}{examplebordercolor}, arc=.5mm
}
\tcolorboxenvironment{remark}{
    enhanced jigsaw, pad at break*=1mm, breakable,
    left=4mm, right=4mm, top=1mm, bottom=1mm,
    colback=white, boxrule=0pt, frame hidden,
    borderline west={0.5mm}{0mm}{examplebordercolor}, arc=.5mm
}


%%%%%%%%%%%%%%%%%%%%%%%%%
% Property Environments %
%%%%%%%%%%%%%%%%%%%%%%%%%

\theoremstyle{property}
\newtheorem{property}{Property}
\newtheorem{prop}{Proposition}

\tcolorboxenvironment{property}{
    enhanced jigsaw, pad at break*=1mm, breakable,
    left=4mm, right=4mm, top=1mm, bottom=1mm,
    colback=propertybackgroundcolor, boxrule=0pt, frame hidden,
    borderline west={0.5mm}{0mm}{propertybordercolor}, arc=.5mm
}
\tcolorboxenvironment{prop}{
    enhanced jigsaw, pad at break*=1mm, breakable,
    left=4mm, right=4mm, top=1mm, bottom=1mm,
    colback=propertybackgroundcolor, boxrule=0pt, frame hidden,
    borderline west={0.5mm}{0mm}{propertybordercolor}, arc=.5mm
}


%%%%%%%%%%%%
% Question %
%%%%%%%%%%%%

\mdfsetup{skipabove=\topskip,skipbelow=\topskip}

\mdtheorem[style=question]{question}{Problem}

\mdfdefinestyle{question}{
  linecolor=black,linewidth=1pt,
  frametitlerule=true,
  frametitlebackgroundcolor=gray!20,
  innertopmargin=\topskip,
}


%%%%%%%%%%%%
% Solution %
%%%%%%%%%%%%

\theoremstyle{solution}
\newtheorem{solution}{Solution}

\tcolorboxenvironment{solution}{
    enhanced jigsaw, pad at break*=1mm, breakable,
    left=4mm, right=4mm, top=1mm, bottom=1mm,
    colback=white, boxrule=0pt, frame hidden,
    borderline west={0.5mm}{0mm}{solutionbordercolor}, arc=.5mm
}


%%%%%%%%%
% Proof %
%%%%%%%%%

% These patches must be placed after \tcolorboxenvironment !
\AddToHook{env/theorem/after}{\colorlet{proofcolor}{theorembordercolor}}
\AddToHook{env/example/after}{\colorlet{proofcolor}{examplebordercolor}}
\AddToHook{env/definition/after}{\colorlet{proofcolor}{definitionbordercolor}}
\AddToHook{env/property/after}{\colorlet{proofcolor}{propertybordercolor}}

\newtheoremstyle{proof}
    {0pt}{0pt}{\normalfont}{0pt}
    {}{\;}{0.25em}
    {{\sffamily\bfseries\color{proofcolor}\thmname{#1}.}
        \thmnote{\normalfont\color{black}~(\textit{#3})}}

\theoremstyle{proof}
\newtheorem*{proofs}{Proof}

\tcolorboxenvironment{proofs}{
    enhanced jigsaw, pad at break*=1mm, breakable,
    left=4mm, right=4mm, top=1mm, bottom=1mm,
    colback=white, boxrule=0pt, frame hidden,
    borderline west={0.5mm}{0mm}{proofcolor}, arc=.5mm
}

\renewcommand{\qedsymbol}{Q.E.D.}
\let\qedsymbolMyOriginal\qedsymbol
\renewcommand{\qedsymbol}{%
    \color{proofcolor}\qedsymbolMyOriginal%
}

\newenvironment{info}{\begin{tcolorbox}[
    arc=0mm,
    colback=white,
    colframe=gray,
    title=Info,
    fonttitle=\sffamily,
    breakable
]}{\end{tcolorbox}}
\newenvironment{terminology}{\begin{tcolorbox}[
    arc=0mm,
    colback=white,
    colframe=green!60!black,
    title=Terminology,
    fonttitle=\sffamily,
    breakable
]}{\end{tcolorbox}}
\newenvironment{warning}{\begin{tcolorbox}[
    arc=0mm,
    colback=white,
    colframe=red,
    title=Warning,
    fonttitle=\sffamily,
    breakable
]}{\end{tcolorbox}}
\newenvironment{caution}{\begin{tcolorbox}[
    arc=0mm,
    colback=white,
    colframe=yellow,
    title=Caution,
    fonttitle=\sffamily,
    breakable
]}{\end{tcolorbox}}
